\section{Das alte Bell10 erscheint aus echten Stromprodukten wieder}
Die tatsächlichen Grad-2-Produkte der 24 Ströme erzeugen einen 50-dimensionalen Raum. Für die 60 Produktvektoren $r_l$ gilt
\[
 \mathcal G=\sum_l\ket{r_l}\bra{r_l}
 =\frac34I_{50}+\frac94P_{10}.
\]
Der alte Bell10-Raum ist damit der eindeutige oberste \emph{Antwortbereich}: zehn Eigenwerte sind $3$, vierzig sind $3/4$. $\mathcal G$ ist kein Hamiltonoperator; die Differenz $9/4$ ist keine Energielücke. \src{21}

Die Level-1-Vierstromform enthält einen echten Klammerterm zusätzlich zu den Wick-artigen Paarungen. Deshalb können einzelne Stromquadrate verschwinden, während für eine vervollständigte Richtung
\[
 \|J_{-1}(w_{l,s})^2\ket\Omega\|^2=\frac59
\]
gilt. Nach Projektion auf Bell10 beträgt das unnormierte Gewicht $5/18$. Eine explizite positive Frequenzabsorption rekonstruiert die alten Effekte auf einem festgelegten Ein-Strom-Bereich. Eine andere Instrumentierung besitzt jedoch dieselben Effekte und andere Nachzustände. Das physische Mess- und Auffrischgesetz bleibt somit zusätzliche Information. \src{21}

\section{Die tatsächliche Coxeterclock erzwingt mehr als D8}
Im festen CAR-Wörterbuch schickt $C$ die 112 D8-Wurzeln in 56 D8- und 56 Spinorwurzeln. Die kumulative $C$-Bahn erreicht bei Potenz acht alle 240 Wurzelrichtungen. Schon D8 zusammen mit einem geeigneten Spinorstrom und seinem Adjungierten schließt durch nichtverschwindende Klammern in den Schritten $114\to170\to240$. \src{22}

Falls also $C$ auf demselben lokalen Stromraum physisch implementiert ist, kann eine rein D8-beschränkte Theorie nicht unverändert bestehen bleiben. Das ist eine notwendige Schließungsstruktur. Es ist noch kein Bau der 128 fehlenden Spinorstromoperatoren auf dem Rohquellenraum. Auch die gewählte Spinorchiralität wird nicht unabhängig selektiert, wenn $C$ schon aus diesem E8-System stammt.

\section{Der feste komplexe Rahmen ist zu klein}
Die 60 alten Strahlen sind viergliedrige $J$-Orbits vollständiger Wurzeln. Die tatsächliche Coxeterwirkung respektiert diese Quotienten bei festem $J$ nicht: Ein Orbit kann auf zwei Zielorbits aufgeteilt werden. Der korrekte kovariante Ausdruck ist
\[
 ([\alpha]_J,J)\longmapsto
 ([C\alpha]_{CJC^{-1}},CJC^{-1}).
\]
Der letzte Vertrag macht diesen Befund quantitativ. Die Abbildung $A=I+J$ erfüllt $A^TA=2I$ und bildet Standard-E8-Wurzeln der Normquadrate zwei auf die tatsächlichen Gaußschen Quellenwurzeln der Normquadrate vier ab. Sie ist eine skalierte Koordinatenbrücke, kein normgleicher Wurzelautomorphismus. \src{29}

Die spezifizierte Koordinatenbrücke transportiert $C$ zu $C_G=ACA^{-1}$ in den Gaußkoordinaten. Sie wird nicht als einzig mögliche Brückenwahl behauptet. Die Rahmenbahn in diesen Koordinaten
\[
 J_r=C_G^rJC_G^{-r}
\]
hat exakt 15 verschiedene Glieder und entsprechende Ereignisalphabete. Die Projektoren transportieren sich kovariant,
\[
 C_G\Pi_\alpha^{J_r}C_G^{-1}=\Pi_{C_G\alpha}^{J_{r+1}}.
\]
Unter $1\le k\le30$ kommutiert $C_G^k$ nur für $k=15,30$ mit $J$; keine Potenz antikommutiert. Eine zusätzliche Aussage im integralen Standard-Gitterrahmen schließt jede gittererhaltende komplexe Struktur $J_*$ mit $J_*^2=-I$ aus, die mit $C$ kommutiert. Im benutzten Ganzzahlrahmen hat $C$ modulo 31 acht verschiedene Eigenwerte; ein kommutierender Operator müsste auf jeder Eigenlinie skalar wirken, obwohl $-1$ in $\F_{31}$ kein Quadrat ist. Antikommutation scheitert bereits an der ungeraden Halbperiode $C^{15}=-I$. \src{29}

Dieser Satz betrifft den integralen Geltungsbereich, nicht beliebige reelle komplexe Strukturen. Zentralisatoren von Coxeterelementen sind zudem bekannte mathematische Objekte; die Sitzung beansprucht nicht die allgemeine Theorie als Neuerfindung. \cite{coxeter}

Positiv ist eine kovariante Familie $H[J_r]$ möglich. Ihre Existenz wählt aber keine Übergänge zwischen den Rahmen, keine Rahmenenergie und keine physische Eichinterpretation. Die Zahl 15 bezeichnet hier die Größe einer Rahmenbahn, keine neue Zahl von Familien oder Raumdimensionen.

\section{Warum neutrale Paketfluktuationen keine acht Cartanströme tragen}
Auf den Ring-Paketzuordnungen wirkt die Quellen-$J$-Transformation als gemeinsamer Skalar. Sie bleibt nach Grammatrix-Orthonormierung skalar. Auf den tatsächlichen acht reellen Cartanrichtungen gilt dagegen $J^2=-I$, also gibt es keinen Fixvektor. Für eine äquivariante lineare Abbildung $T$ müsste gelten
\[
 T(Jv)=\operatorname{Ad}_{U_J}T(v)=T(v),
\]
und damit $T=0$. Eine mit der Symmetrie verschränkte Dressierung ändert diesen Ausschluss innerhalb des neutralen Raums nicht. \src{29}

Die richtige Lehre lautet nicht, neutrale Observablen zu verwerfen. Ihre Produktstruktur kann über geladene Zwischenzustände vermittelt sein:
\[
 P[A,B]P=[PAP,PBP]+PAQBP-PBQAP,\qquad Q=I-P.
\]
Eine bloße Kompression verliert die letzten beiden Terme. Für die gesuchte Stromantwort müssen deshalb die geladenen Komplementwege mit Energie und spektralem Gewicht mitgeführt werden. Eine passende Anzahl neutraler Zustände reicht nicht.

\section{Was die Clock an der Zeitregel bereits erzwingt}
Unter der ausdrücklichen Prämisse einer phasendiagonalen affinen Stromwirkung besitzt ein Generator die Form
\[
 [H,J_n^\alpha]=(\alpha(h)-an)J_n^\alpha.
\]
Die tatsächliche Coxeterwirkung mit charakteristischem Polynom $\Phi_{30}$ erzwingt $h=0$. Auf dem gemeinsamen irreduziblen Stromraum bleibt somit
\[
 H=aL_0+c.
\]
Die Regel ist bedingt stark: Ist der gemeinsame affine Prozess erst vorhanden, bleibt nur eine positive Zeitkalibrierung. Sie beweist jedoch nicht die bislang vorausgesetzte Invarianz des linearen Stromraums. \src{28}

Eine diskrete Clock und ein gemeinsames Vakuum allein genügen nicht. Gegenbeispiele der Form $H_b=L_0+4bL_0(L_0-1)$ können den geprüften Viertelschritt und den ersten Grad bewahren, aber höhere Grade anders bewegen. Dabei erfüllt jeder Kommutator $\operatorname{ad}_{H_b}$ weiterhin die Leibnizregel. Was scheitert, ist die state-unabhängige skalare Frequenz auf dem linearen Stromraum, nicht die Leibnizregel selbst. \src{28}

\section{Der positive Fermi-Patch: ein kontrollierter Anfang}
Der letzte Herkunftstest arbeitet wieder mit der tatsächlichen CAR-Quelle. Die dortige Einteilchendispersion ist $\varepsilon(k)=-2\cos k$. Auf einem festen, nicht umlaufenden Patch um eine Fermistelle liefert die Skalierung
\[
 e_L(r)=\frac L{2\pi}\sin\frac{2\pi r}L
\]
die Abschätzung
\[
 \bigl|e_L(r)-e_L(r+n)+n\bigr|
 \le\frac{(2\pi)^2}{6L^2}\bigl(|r|^3+|r+n|^3\bigr).
\]
Für passend verschachtelte Träger auf dem festen zusammenhängenden Halbzahl-Patch gilt für die positiven Moden $1,2,3$ zugleich der echte Modenkommutator
\[
 [E^{12}_1,E^{23}_2]=E^{13}_3.
\]
Unter der zusätzlichen gleichen-Spezies-Kinetik $h_8=I_8\otimes h_{\mathrm{v454}}$ gehen die Frequenzen mit kontrolliertem Fehler gegen $-1,-2,-3$; am anderen Fermiast wechseln ihre Vorzeichen. Das ist mehr als eine bloß passende Zustandszählung: Komposition und Zeitfrequenz werden auf demselben ausgeschnittenen CAR-Prozess gemeinsam geprüft. \src{29}

Die Grenze bleibt entscheidend. Acht identische Kinetiken entstehen nur nach der zusätzlichen Wahl $h_8=I_8\otimes h_{\mathrm{v454}}$. Acht Slots erzwingen diese Wahl nicht. Ein fester Patch liefert noch keinen Adjungiertenabschluss, keinen gefüllten Seezustand, keinen Level-1-Zentralterm und keine gemeinsame GNS-Domäne. Die 128 Spinorstromoperatoren fehlen auf demselben Rohraum weiterhin.

Rigorose Gitter-Kontinuum-Konstruktionen für freie Fermionen und WZW-Modelle bieten einen passenden Vergleichsrahmen für die nächsten Konvergenzschritte. Ihre Voraussetzungen müssen aber am TFPT-Quellobjekt nachgewiesen werden. Sie ersetzen weder die Auswahl der Kinetik noch die P1/P2-Identifikation. \cite{osborne}

\chapter{Was die Sitzung insgesamt verändert hat}
\section{Die positiven Ergebnisse zusammen lesen}
Die Sitzung hat die Ursprungsidee nicht auf eine weitere Zahlenübereinstimmung reduziert. Sie hat mehrere belastbare Verbindungen hergestellt:
\begin{enumerate}
\item \textbf{Das Alphabet besitzt präzise Komposition.} Integrale Trialität, native CAR-Wörterbücher, Chevalleyzeichen, quartischer Übertrag und Rahmenwirkungen sind explizit. Falsche Gleichsetzungen werden durch die tatsächlichen Wirkungen ausgeschlossen.
\item \textbf{Ein gemeinsamer Ereignisprozess kann verschiedene Auslesungen tragen.} Die volle 60er-Quelle und der Petersen-Clock sind mit einem phasentreuen Record zusammengeführt. Die kontinuierliche und affine Fortsetzung ist jeweils mit ihren Zusatzannahmen benannt.
\item \textbf{Die Quelle kann Teil des Quantenzustands sein.} Das Einzelpaket und mehrere überlappende endliche Systeme besitzen kontrollierte Spektralsätze. Die lokale Reinheit muss dabei nicht auf alle überlappenden Teile übertragen werden.
\item \textbf{Der richtige Detektor kann eine übersehene Bewegung sichtbar machen.} Phasen- und Bindungsoperatoren sehen einen niedrigen gleichphasigen Partner, den lineare Wurzelkoordinaten aufgrund von Symmetrie nicht sehen können.
\item \textbf{Ein erster Prozessanschluss an Stromfrequenzen ist kontrolliert.} Der Fermi-Patch liefert einen nichttrivialen Modenkommutator und passende Frequenzen mit einer expliziten $L^{-2}$-Schranke. Er ist ein Anfang einer gemeinsamen Grenzwertkonstruktion, kein vollständiges Stromnetz.
\end{enumerate}

Die zusammenhängende Lehre ist: Information verschwindet oft nicht aus dem System, sondern nur aus der gerade gewählten Auslesung. Für eine Ursprungstheorie muss die kleinere Beschreibung nicht bloß heutige Werte, sondern die relevanten möglichen Fortsetzungen bewahren.

\section{Die wichtigsten Korrekturen im Verlauf}
\begin{longtable}{@{}P{.25\textwidth}P{.42\textwidth}P{.25\textwidth}@{}}
\toprule
Frühere naheliegende Deutung & Heutiger belastbarer Stand & Konsequenz\\\midrule
\endfirsthead
\toprule Frühere Deutung & Heutiger Stand & Konsequenz\\\midrule
\endhead
Zustand bzw. Vorzeichenprojektor bestimmt den Generator & Energiegrößen und Kohärenzdynamik bleiben im Allgemeinen frei. & Mehrzeitgesetz oder Generatorursprung nötig.\\[4pt]
Acht reale Richtungen sind die acht nativen CAR-Moden & Charaktere und Skalarkörper unterscheiden sich. & Geladenes Operatorwörterbuch statt Dimensionsvergleich.\\[4pt]
Gruppenordnung 1920 identifiziert $W(D_5)$ & Der Ereignisrahmen ist $\mathrm{ASL}_2(4):2$ mit anderer Wirkung. & Den vollen phasentreuen Rahmen speichern.\\[4pt]
Gleicher Ereignisendpunkt liefert dieselbe Zeit & Unterschiedliche positive Generatoren und Refreshregeln sind möglich. & Zwischenbewegung und Mehrzeitprozess prüfen.\\[4pt]
Kleine Quellenlücke bedeutet masseloses Teilchen & Im starken Zweig ist sie kollektives Ausrichtungstunneln. & Lokale Antwort und Dispersion über dem Vakuum nötig.\\[4pt]
Tiefer Paarsektor ist das Vollvakuum & Explizite erlaubte Zustände liegen darunter. & Alle konkurrierenden Quellsektoren berücksichtigen.\\[4pt]
Zweiquellendublett bereits ursprünglich vollständig bewiesen & In .24 Annahme G; erst .26 exakter Vollraumsatz. & Historische Numerik und späteren Beweis getrennt zitieren.\\[4pt]
Dreiquellenresonanz liegt nachgewiesen über dem Vollvakuum & Invarianter Zweiblock exakt; globaler Komplementboden offen. & Antworten vorerst an diesen Block binden.\\[4pt]
Vierquellensektor $K=1$ ist das ganze System & Zehn $K=3$-Masken erst in .28 ausgeschlossen. & Vollsatz erst ab .28.\\[4pt]
Ising-Kompression beweist kritische Materie & .29 liefert einen tieferen Vollraum-Gegenzeugen. & Nackte Paketroute ausgeschlossen; Dressing offen.\\[4pt]
Bewegliche Bindung bedeutet bewegliche Geometrie & Graphinzidenz bleibt ohne zusätzliche Übergänge erhalten. & Ortsalgebren und Inzidenzprozess herleiten.\\[4pt]
Gleiche Effekte transportieren den ganzen Prozess & Nachzustände, geladene Zwischenräume und Zeit können verschieden sein. & Instrumente und Mehrzeitkorrelationen vergleichen.\\\bottomrule
\end{longtable}

\section{Warum wir noch keine Gesamtlösung haben}
Der Engpass ist nicht, dass uns eine weitere besonders elegante Zahl fehlt. Er besteht darin, dass die tragenden positiven Resultate noch unter verschiedenen Zusatzvoraussetzungen stehen. Der eine Satz setzt einen Quellen-Hamiltonoperator voraus; der nächste eine affine Vakuumdarstellung; ein dritter eine Markierung und einen Record; ein vierter acht gleiche freie Kinetiken. Diese Voraussetzungen werden nicht dadurch gemeinsam wahr, dass ihre Ausgaben gut zusammenpassen.

Eine vollständige Lösung müsste diese Auswahl selbst leisten und anschließend zeigen, dass im selben Zweig die richtigen lokalen geladenen Felder, Raumzeitantworten und physikalischen Kopplungen bestehen. Die bereits vorhandene TFPT-Architektur gibt dafür viele strenge Rückbedingungen. Sie ist deshalb mehr als eine Sammlung statischer Zahlen. Die Rückbedingungen ersetzen aber nicht den gemeinsamen Existenz- und Herkunftsnachweis.

\noindent\begin{tabularx}{\textwidth}{@{}lYY@{}}
\toprule
Vertrag & Benötigter Gesamtabschluss & Beitrag dieser Sitzung\\\midrule
T1 & Ursprung von P1/P2 und Dimension & Zusätzliche Geometrie-, Quellen- und Zeitwahlen isoliert; nicht hergeleitet.\\
T2 & Tatsächliche E8-Randtheorie und Grenzwert & Native Wörterbücher, Produkte, Clockschließung und bedingter Fermi-Patch.\\
T3 & Gemeinsamer lokaler unitärer 3+1D-Parent & Endliche lokale Modelle und Gegenbeispiele; kein gemeinsamer Parent.\\
T4 & Chirales SM, kontrollierte Spiegelentkopplung & Ladungs- und Intertwinerhindernisse präzisiert; keine vollständige chirale Feldrealisierung.\\
T5 & Wechselwirkendes Kontinuum, Lorentz-/IR-Grenze & Falsche Gap- und Paketkritikalitätsidentifikationen ausgeschlossen; kontrollierter Teilpatch.\\
T6 & Kopplungen und Neutrinodynamik & Freie $J,\kappa,\mu$, Geschwindigkeiten und Instrumentdaten sichtbar; keine Auswahl.\\
T7 & Universeller masseloser Quantenspin-2-Sektor & Keine solche Ableitung; interne Acht und kleine Lücken ausdrücklich getrennt.\\
T8 & Anfangszustand und gemeinsame Erzeugungsfunktion & Energieauswahl von Präparation unterschieden; Prozessgedächtnis rekonstruiert.\\\bottomrule
\end{tabularx}

\chapter{Die größten anschließenden Forschungsfragen}
\section{Priorität 1: dieselbe Quelle, dieselben geladenen Operatoren}
\textbf{Frage.} Lässt sich aus P1/P2 eine einzige lokale Realisierung gewinnen, in der die tatsächlichen 16 Majoranas, die 120 D8-Ströme, die 128 Spinorströme und die markierten $C,J$-Wirkungen gemeinsam handeln?

\textbf{Warum dies den größten Hebel hat.} Der zentrale Engpass ist die erste operatorielle Verbindung zwischen Rohquelle und der bereits sehr präzisen E8-Stromtheorie. Ohne sie können Aussagen über Quelle, Vakuum und Zeit weiterhin aus verschiedenen Hilfsmodellen stammen. Die neutrale Paketspanne reicht nach dem zentralen Auswahltest nicht aus.

\textbf{Nächster entscheidender Nachweis.} Auf einem expliziten gemeinsamen dichten Kern mindestens einen geladenen Spinoroperator, seinen Adjungierten und einen nichtverschwindenden D8-Spinor-Klammerpfad konstruieren. Die Abbildung muss Adjunktion, Kokzykluszeichen, Ladungen und die \emph{tatsächliche} Clockwirkung erhalten. Ihre Herkunft aus P1/P2 muss die erste zusätzliche Wahl offenlegen.

\textbf{Erfolg.} Ein nichtnull lokaler geladener Operator mit korrekter Antwortnorm, Klammer und Clocktransformation im selben Vakuumprozess. Danach darf der vollständige 128er-Abschluss bearbeitet werden.

\textbf{Abbruch- oder Umkehrkriterium.} Scheitert bereits die zentrale Ladung oder ein Charaktertest, ist diese konkrete Identifikation zu verwerfen. Eine bloße Erweiterung der neutralen Matrixdimension repariert sie nicht.

\section{Priorität 2: aus einer Clock eine physische Zeit gewinnen}
\textbf{Frage.} Warum wählt die Rohquelle gerade einen gemeinsamen Generator $H=aL_0+c$ und, soweit erforderlich, acht gleiche Kinetiken?

Der bedingte Clockklassifikationssatz und der Fermi-Patch zeigen ein klares Ziel. Noch offen sind Astwahl, gemeinsame Geschwindigkeit, Vakuumfüllung, Normalordnung und die Grenzwertdomäne. Der diskrete Clock-Schritt allein ist nach den Gegenbeispielen zu schwach.

\textbf{Nächster Nachweis.} Den vorhandenen positiven Patch zu einer unter Adjunktion und Kommutatoren geschlossenen Folge ausbauen. Für zwei nichtkommutierende Moden und ihre Klammer dieselbe Skalierung $a$ kontrollieren; anschließend den zentralen Term aus dem tatsächlichen Seezustand gewinnen. Die Erweiterung des Patches muss mit einer Fehlerabschätzung erfolgen, statt den festen-Patch-Satz stillschweigend auf alle Moden zu übertragen.

\textbf{Erfolg.} Gemeinsame starke oder geeignete korrelationsbasierte Konvergenz auf einer GNS-Domäne, positives Vakuum, Level-1-Zentralterm und keine separat eingesetzten Geschwindigkeiten. Die verbleibende metrologische Zeitkalibrierung darf benannt extern sein; relative Kinetiken dürfen nicht frei bleiben.

\textbf{Abbruchkriterium.} Ein persistenter unterschiedlicher Frequenzfaktor auf einem Klammerpfad widerlegt die angenommene gemeinsame Zeitrealisierung.

\section{Priorität 3: eine dressierte kritische Phase über dem echten Vakuum}
\textbf{Frage.} Gibt es in der unveränderten Quellenfamilie bei festen positiven Kopplungen einen großen Zustandszweig mit lokalen $z=1$-Anregungen und richtigem geladenen Antwortgewicht?

Die nackte Ising-Paketfamilie ist ausgeschlossen. Offen bleibt eine stärker gemischte Phase. Sie ist nur dann zielführend, wenn ein plausibler Herkunftsanschluss der Hamiltonfamilie erhalten bleibt; eine beliebige Reparaturwechselwirkung wäre keine Lösung der ursprünglichen Frage.

\textbf{Nächster Nachweis.} Ein quasilokales Dressing oder eine Feshbach-/Schur-Elimination mit kontrolliertem Rest konstruieren. Der Rest muss klein gegenüber der \emph{mit der Größe schrumpfenden} relevanten Lücke sein. Gleichzeitig ist die Energie unter allen vorhandenen Gegenzeugen zu halten und mindestens ein korrekt markierter geladener Spektralzweig zu verfolgen.

\textbf{Erfolg.} Zusammengehörig nachgewiesen werden müssen (i) das Vollraumvakuum beziehungsweise eine kontrollierte thermodynamische Phase, (ii) lineare Dispersion lokaler Antworten und (iii) nichtverschwindendes skaliertes Stromgewicht. Ein alleiniger Fit von $N\Delta_N$ genügt nicht.

\textbf{Abbruchkriterium.} Extensiver nichtkontrollierter Rest, ein tieferer Konkurrenzzustand oder verschwindendes geladenes Gewicht schließen den jeweiligen Kandidaten aus. Die bereits ausgeschlossene nackte Familie wird ohne neue begründete Voraussetzung nicht erneut optimiert.

\section{Priorität 4: ein physisches Instrument samt Gedächtnis}
\textbf{Frage.} Welche aus der Quelle bestimmte Operation entscheidet nicht nur die Messwahrscheinlichkeit, sondern auch Nachzustand, Record, Ereigniszeit und Auffrischung?

Die Sitzung hat wiederholt Effekte und Zustandsisometrien rekonstruiert, ohne damit den ganzen Prozess zu bestimmen. Der Unterschied zwischen gehaltener und aufgefrischter Quelle ist ein direkter Zwei-Zeit-Zeuge. Er macht diese Frage experimentell beziehungsweise operational formulierbar.

\textbf{Nächster Nachweis.} Auf dem gemeinsamen Quell-/Stromraum eine CP-Instrumentfamilie samt Record angeben und die zwei aufeinanderfolgenden nativen Ereignisse vergleichen. Die erhaltene Beschreibung muss dieselben Nachzustände und Mehrzeitkorrelationen liefern, nicht nur denselben POVM-Effekt. Markierung und Rahmen müssen auch zwischen Ereignissen kovariant behandelt werden.

\textbf{Erfolg.} Ein quellengewähltes Mehrzeitgesetz, das die Petersen-Auslese, die Bell10-Antwort und die geladene Stromantwort gleichzeitig reproduziert. \textbf{Abbruchkriterium.} Ein einziger zulässiger Folgeeingriff, dessen Antwort die komprimierte Beschreibung nicht bestimmen kann, widerlegt deren behauptete operative Vollständigkeit.

\section{Priorität 5: aus Beziehungen wirkliche Lokalität rekonstruieren}
\textbf{Frage.} Welche voneinander unterscheidbaren lokalen Algebren und welche Dynamik ihrer Inzidenz ergeben sich aus dem gemeinsamen Prozess, ohne die gesuchte Raumgeometrie schon einzusetzen?

Die Sitzung erlaubt die Hypothese, dass Bindungen und ihre Rekonfigurationen grundlegender sind als feste Teilchenorte. Sie beweist diese Hypothese nicht. Die blockdiagonale Grapherweiterung erzeugt keine Kantenbewegung und bevorzugt ohne neue Regel Matchings statt zusammenhängender Räume.

\textbf{Nächster Nachweis.} Zunächst zwei wirklich verschiedene kommutierende lokale Algebren rekonstruieren, die unter der tatsächlichen Dynamik unterscheidbare Antworten behalten. Danach ein aus der Quelle abgeleitetes Matrixelement zwischen Inzidenzmustern oder eine gleichwertige operationale Ortsrekonstruktion vorlegen. Relative Graphenergien dürfen nicht unbemerkt frei gewählt werden.

\textbf{Erfolg.} Eine stabile zusammenhängende Phase mit kontrollierter Ausbreitung und operational gemessener Dimension. Eine spektrale Zahl nahe drei oder eine gezeichnete 3D-Einbettung allein genügt nicht. \textbf{Abbruchkriterium.} Werden verschiedene Orte durch die maßgebliche Algebra identifiziert, trägt diese Algebra die behauptete Geometrie nicht.

\section{Priorität 6: chirale Materie, 3+1D und Gravitation gemeinsam}
\textbf{Frage.} Kann dieselbe erfolgreiche Realisierung anschließend das chirale Standardmodell und einen universell gekoppelten masselosen Spin-2-Zweig tragen?

Dies ist die größte Endfrage, aber nicht der kleinste nächste Rechenschritt. Interne Spinordarstellungen sind keine Raumzeitspinoren. Ein eindimensionaler linearer Zweig ist keine 3+1-dimensionale chirale Theorie. Anomaliefreiheit allein beweist weder Spiegelentkopplung noch die erforderlichen Wechselwirkungen.

\textbf{Erfolgskriterien.} Ein gemeinsamer lokaler Parent muss Ladungen und Graduierung, chirale Felder, kontrollierte Spiegelbehandlung, Kontinuum, Kopplungsverhältnisse und universelle gravitative Antwort im selben Zustandszweig tragen. Dazu gehört ein begründeter Anfangszustand oder ein Präparationsprinzip. Die früheren TFPT-Auslesungen sind dann aus diesem Prozess zurückzugewinnen.

\textbf{Arbeitsreihenfolge.} Zuerst die engste tragende Herkunfts- und Operatorfrage aus Priorität 1 schließen; Zeit und Instrument daran testen; große Phase und Lokalität nur mit diesem gemeinsamen Bezug entwickeln. Die sechs Fragen sind keine Einladung zu sechs beliebig neuen Modellwelten. Sie sind aufeinander angewiesene Nachweise für dasselbe Ziel.

\section{Eine notwendige Bestandsbereinigung}
Unabhängig von diesen großen Fragen bleibt ein kleinerer dokumentarischer Nachweis sinnvoll: Der globale Drei-Quellen-Komplementboden aus dem früheren Bericht sollte entweder mit vollständigem Zertifikat rekonstruiert oder die entsprechende Vollraumbehauptung dauerhaft auf den invarianten Bindungsblock begrenzt werden. Gleiches gilt für die unreplayte allgemeine Leakageformel. Das erhöht die Verlässlichkeit des Bestands; es ist nicht mit einem Fortschritt zur vollständigen Weltphysik zu verwechseln.

\chapter{Provenienz, Reproduzierbarkeit und Reichweite}
\section{Wie dieser Bericht erstellt wurde}
Dieser Bericht ist eine neue redaktionelle Synthese der Arbeitssitzung. Die Originalbeweise wurden gezielt mit ihren Vertragsindizes abgeglichen. Frühere Formulierungen aus dem Gespräch wurden bei Konflikten durch den präziseren späteren Status ersetzt. Neue Quellenmodelle, neue Naturkonstanten oder neue Beweissätze wurden für diese Dokumentation nicht eingeführt.

Der Repository-HEAD beim Quellenabgleich war
\begin{center}\small\ttfamily 66b91e40e245569f06ab440ead80f446c9be0ee5.\end{center}
Die Sitzungsergebnisse liegen zusätzlich im Arbeitsbaum. Der Commit allein identifiziert deshalb nicht den dokumentierten Forschungsstand. Das begleitende JSON-Manifest bindet die verwendeten Dateien mit SHA-256. Die Quellenregistereinträge verweisen auf Beweisdateien und Vertragsindizes; deren Zertifikate und Replaydaten bleiben die ausführlichen Prüfobjekte.

Der Theoriegraph hatte beim Abgleich 5914 Knoten und 63\,758 Kanten. Er wurde zur Zuordnung benutzt, nicht als Beweis. Die 29 nummerierten Ergebnispositionen verteilen sich auf 28 Vertragsindizes: .04 und .05 teilen einen Vertrag; .10 heißt \texttt{UR.SEAM.LOCALIZED\_SECTOR.10}. Die Gesamtverdicts dieser Forschungsverträge bleiben \texttt{PARTIAL}.

\section{Was im Dokumentationslauf wirklich geprüft wurde}
Der Repository-Audit \texttt{bash build.sh audit} endete am 20.09.2026 mit Exitcode null. Er meldete unter anderem 266 bestandene RH-Prüfungen, einen konsistenten Katalog und einen konsistenten Graphstand. Das ist ein Integritätsbefund des angegebenen Auditmodus, keine neue physische Theorievalidierung. Der schnelle Modus übersprang Module; fünf erklärte Lean-\texttt{sorry}-Warnungen, fünf Graphzykluswarnungen und zwei erwartete Drifts ungepinnter lebender Dokumente bleiben ausdrücklich sichtbar.

Die großen Spektralrechnungen dieser Sitzung wurden zur Erstellung des Dossiers nicht sämtlich wiederholt. Numerische Zahlen werden daher als dokumentierte Ergebnisse und nicht als frischer Vollraumlauf bezeichnet. Die Quellenhashes wurden abgeglichen. LaTeX-Kompilation, interne Verweise und die gerenderte Seitendarstellung des neuen PDFs werden gesondert kontrolliert. Ein sauber gesetztes PDF ersetzt weder eine unabhängige mathematische Begutachtung noch einen Beweisassistenten.

Die LaTeX-Quelle ist eigenständig; Diagramme werden darin erzeugt. Zur erneuten PDF-Erstellung genügt in einer TeX-Live-Umgebung \texttt{latexmk -pdf} mit dem Namen dieser \texttt{.tex}-Datei. Die begleitende \texttt{.sources.json}-Datei dokumentiert Quellenstand und Ausgabehashes.

\section{Notation und typische Verwechslungen}
\begin{longtable}{@{}P{.2\textwidth}P{.72\textwidth}@{}}
\toprule Zeichen & Bedeutung\\\midrule\endfirsthead
\toprule Zeichen & Bedeutung\\\midrule\endhead
$J$ in $H$ & Nichtnegative Ereigniskopplung; eine Zahl.\\
$J,J_G$ als Matrix & Tatsächliche Viertelclock bzw. Gaußsche komplexe Struktur; im jeweiligen Quellenrahmen ausdrücklich angegeben. Kein Kopplungswert.\\
$J_{\F_4}$ & Ordnung-3-Struktur des binären Rahmenkerns; erfüllt $J_{\F_4}^2+J_{\F_4}+I=0$.\\
$J_n^a$ & Affiner Strommodus; weder Kopplung noch Viertelmatrix.\\
$C$ & Tatsächliche Coxeterwirkung auf dem markierten E8-System.\\
$\mathcal C_{60}$ & Quantenkanal aus 60 Quellenereignissen; kein Coxeterelement.\\
$G$ & Je nach explizitem Abschnitt Klebeclock, Grammatrix oder gemeinsame Ereignisinvolution $r\otimes r$; Antwortoperator in Kapitel 8 als $\mathcal G$ geschrieben.\\
$\kappa$ in $H$ & Quellenbewegungskopplung. Die invariant-bilineare Form $\kappa(a,b)$ in Stromformeln ist ein anderes, durch Argumente erkennbares Objekt.\\
$m,N$ & $m$ Quellen einer offenen Kette mit $m+1$ Registern; $N$ Quellen und Register auf dem Ring.\\
$\Delta$ & Energielücke nur bei einem benannten Hamiltonoperator. Gram- und Antwortlücken sind keine Energien.\\
$P,Q=I-P$ & Jeweils ausdrücklich gewählter Teilraum und Komplement. Ein $PHP$ ist ohne Invarianz oder Fehlerkontrolle keine effektive Volltheorie.\\\bottomrule
\end{longtable}

\appendix

\chapter{Quellen- und Ergebnisregister}
Alle Kennungen S01--S29 verweisen auf die zugehörigen Originalverträge. Maßgeblich sind deren konkrete Teilverdicts und Voraussetzungen; das gemeinsame Gesamtverdict lautet jeweils \texttt{PARTIAL}. Ein Link auf eine Beweisdatei ist keine Behauptung unabhängiger Begutachtung. Die vollständigen SHA-256-Werte aller hier gebundenen Dateien stehen im begleitenden Manifest.

\textbf{Pfadwurzeln.} \texttt{TFPT/} bezeichnet das Repository \path{/Users/stefanhamann/Projekte/tfpt-theoryv4/}. \texttt{SESSION/} bezeichnet \path{/Users/stefanhamann/Documents/Codex/2026-09-18/unte-2/}. Diese Präfixe machen die Fundstellen lesbar; das JSON-Manifest enthält die absoluten Pfade.

\section{Die nummerierte Ergebnisfolge}
\subsection*{S01 — Integrale Trialität}\label{s:01}
{\small\path{UR.COMPILER.INTEGRAL_TRIALITY.01}\par}
Vertrag: \path{TFPT/experiments/theory-contracts/compiler-integral-triality-20260918/contract_index.json}\par{\footnotesize SHA-256 (Anfang): \texttt{54d6580e9413d4c8}.}\par
Herleitung: \path{TFPT/experiments/theory-contracts/compiler-integral-triality-20260918/PROOF.txt}\par{\footnotesize SHA-256 (Anfang): \texttt{123e076cff29adca}.}\par
\subsection*{S02 — Starrheit der integralen Trialität}\label{s:02}
{\small\path{UR.COMPILER.TRIALITY_RIGIDITY.02}\par}
Vertrag: \path{TFPT/experiments/theory-contracts/compiler-triality-rigidity-20260918/contract_index.json}\par{\footnotesize SHA-256 (Anfang): \texttt{43da16e1aa0762bb}.}\par
Herleitung: \path{TFPT/experiments/theory-contracts/compiler-triality-rigidity-20260918/PROOF.txt}\par{\footnotesize SHA-256 (Anfang): \texttt{8cbc553eb0427d42}.}\par
\subsection*{S03 — Natives Wurzel- und CAR-Wörterbuch}\label{s:03}
{\small\path{UR.COMPILER.NATIVE_ROOT_DICTIONARY.03}\par}
Vertrag: \path{TFPT/experiments/theory-contracts/compiler-native-root-dictionary-20260918/contract_index.json}\par{\footnotesize SHA-256 (Anfang): \texttt{38c3007af98fafff}.}\par
Herleitung: \path{TFPT/experiments/theory-contracts/compiler-native-root-dictionary-20260918/PROOF.txt}\par{\footnotesize SHA-256 (Anfang): \texttt{b5de7459d2c849ea}.}\par
\subsection*{S04 — E6-Auslese und Determinantenergänzung}\label{s:04}
{\small\path{UR.COMPILER.E6_READOUT_CLOSURE.04}\par}
Vertrag: \path{TFPT/experiments/theory-contracts/compiler-e6-seam-readout-20260918/contract_index.json}\par{\footnotesize SHA-256 (Anfang): \texttt{6478c8ffd27d7901}.}\par
Herleitung: \path{TFPT/experiments/theory-contracts/compiler-e6-seam-readout-20260918/PROOF.txt}\par{\footnotesize SHA-256 (Anfang): \texttt{a40b401a71eca7a0}.}\par
\phantomsection\label{s:05}\textbf{S05 — Nahtcharakter auf E8.} \path{UR.COMPILER.SEAM_E8_CHARACTER_EXTENSION.05} ist ein eigenständiges verwandtes Resultat desselben Vertrags und derselben Beweisdatei; kein fehlender fünfter Beweisordner.

\subsection*{S06 — Affine Vakuumantwort und E6-Kubik}\label{s:06}
{\small\path{UR.COMPILER.VACUUM_CURRENT_CUBIC.06}\par}
Vertrag: \path{TFPT/experiments/theory-contracts/compiler-vacuum-current-cubic-20260918/contract_index.json}\par{\footnotesize SHA-256 (Anfang): \texttt{32daea34ae13bd70}.}\par
Herleitung: \path{TFPT/experiments/theory-contracts/compiler-vacuum-current-cubic-20260918/PROOF.txt}\par{\footnotesize SHA-256 (Anfang): \texttt{d65189b806facbb1}.}\par
\subsection*{S07 — Quellenrang und Clockwörterbuch}\label{s:07}
{\small\path{UR.COMPILER.SOURCE_CHANNEL_GATE.07}\par}
Vertrag: \path{TFPT/experiments/theory-contracts/compiler-source-channel-gate-20260918/contract_index.json}\par{\footnotesize SHA-256 (Anfang): \texttt{235cc326022476b1}.}\par
Herleitung: \path{TFPT/experiments/theory-contracts/compiler-source-channel-gate-20260918/PROOF.txt}\par{\footnotesize SHA-256 (Anfang): \texttt{7c9e47e12d06c94c}.}\par
\subsection*{S08 — Operations- und Prozessrekonstruktion}\label{s:08}
{\small\path{UR.COMPILER.PROCESS_RECONSTRUCTION.08}\par}
Vertrag: \path{TFPT/experiments/theory-contracts/compiler-process-reconstruction-20260918/contract_index.json}\par{\footnotesize SHA-256 (Anfang): \texttt{baa1c6d47d1a955b}.}\par
Herleitung: \path{TFPT/experiments/theory-contracts/compiler-process-reconstruction-20260918/PROOF.txt}\par{\footnotesize SHA-256 (Anfang): \texttt{c624193a3fde4c2b}.}\par
\subsection*{S09 — Quartischer Übertrag}\label{s:09}
{\small\path{UR.COMPILER.QUARTIC_LIFT.09}\par}
Vertrag: \path{TFPT/experiments/theory-contracts/compiler-quartic-carry-20260919/contract_index.json}\par{\footnotesize SHA-256 (Anfang): \texttt{0916af9829945938}.}\par
Herleitung: \path{TFPT/experiments/theory-contracts/compiler-quartic-carry-20260919/PROOF.txt}\par{\footnotesize SHA-256 (Anfang): \texttt{e38c31499b54bab6}.}\par
\subsection*{S10 — Lokalisierter Spinorsektor}\label{s:10}
{\small\path{UR.SEAM.LOCALIZED_SECTOR.10}\par}
Vertrag: \path{TFPT/experiments/theory-contracts/source-localized-spinor-20260919/contract_index.json}\par{\footnotesize SHA-256 (Anfang): \texttt{1302407ab999b263}.}\par
Herleitung: \path{TFPT/experiments/theory-contracts/source-localized-spinor-20260919/PROOF.txt}\par{\footnotesize SHA-256 (Anfang): \texttt{2a08b4d68fc27030}.}\par
\subsection*{S11 — Lokaler Stromursprung und neutraler Zellentest}\label{s:11}
{\small\path{UR.COMPILER.LOCAL_CURRENT_ORIGIN.11}\par}
Vertrag: \path{TFPT/experiments/theory-contracts/compiler-local-current-origin-20260919/contract_index.json}\par{\footnotesize SHA-256 (Anfang): \texttt{9cad6bb44b631fc8}.}\par
Herleitung: \path{TFPT/experiments/theory-contracts/compiler-local-current-origin-20260919/PROOF.txt}\par{\footnotesize SHA-256 (Anfang): \texttt{5270d1f86ae6c3d1}.}\par
\subsection*{S12 — Quartischer Phasentransport}\label{s:12}
{\small\path{UR.COMPILER.QUARTIC_HOLONOMY.12}\par}
Vertrag: \path{TFPT/experiments/theory-contracts/compiler-quartic-holonomy-20260919/contract_index.json}\par{\footnotesize SHA-256 (Anfang): \texttt{3bd9449d2f437753}.}\par
Herleitung: \path{TFPT/experiments/theory-contracts/compiler-quartic-holonomy-20260919/PROOF.txt}\par{\footnotesize SHA-256 (Anfang): \texttt{5b5664c9704e9b60}.}\par
\subsection*{S13 — Korrelierte Ereignisse und Petersen-Clock}\label{s:13}
{\small\path{UR.COMPILER.CORRELATED_EVENT_CLOCK.13}\par}
Vertrag: \path{TFPT/experiments/theory-contracts/compiler-correlated-event-clock-20260919/contract_index.json}\par{\footnotesize SHA-256 (Anfang): \texttt{dcae87ea82c76b4f}.}\par
Herleitung: \path{TFPT/experiments/theory-contracts/compiler-correlated-event-clock-20260919/PROOF.txt}\par{\footnotesize SHA-256 (Anfang): \texttt{bf1f85c99f5f89a0}.}\par
\subsection*{S14 — Phasenbewusste Rahmengruppe}\label{s:14}
{\small\path{UR.COMPILER.FRAME_GROUP.14}\par}
Vertrag: \path{TFPT/experiments/theory-contracts/compiler-frame-group-20260919/contract_index.json}\par{\footnotesize SHA-256 (Anfang): \texttt{e70773cc99dd2a37}.}\par
Herleitung: \path{TFPT/experiments/theory-contracts/compiler-frame-group-20260919/PROOF.txt}\par{\footnotesize SHA-256 (Anfang): \texttt{60969a33006b3134}.}\par
\subsection*{S15 — Kontinuierlicher Ereignisursprung}\label{s:15}
{\small\path{UR.COMPILER.CONTINUOUS_EVENT_ORIGIN.15}\par}
Vertrag: \path{TFPT/experiments/theory-contracts/compiler-continuous-event-origin-20260919/contract_index.json}\par{\footnotesize SHA-256 (Anfang): \texttt{0f0f138705e6ed02}.}\par
Herleitung: \path{TFPT/experiments/theory-contracts/compiler-continuous-event-origin-20260919/PROOF.txt}\par{\footnotesize SHA-256 (Anfang): \texttt{545cb24fd942e67d}.}\par
\subsection*{S16 — Affiner Nachfahrenanschluss}\label{s:16}
{\small\path{UR.COMPILER.CURRENT_DESCENDANT.16}\par}
Vertrag: \path{TFPT/experiments/theory-contracts/compiler-current-descendant-20260919/contract_index.json}\par{\footnotesize SHA-256 (Anfang): \texttt{d27ef5b55d8716a6}.}\par
Herleitung: \path{TFPT/experiments/theory-contracts/compiler-current-descendant-20260919/PROOF.txt}\par{\footnotesize SHA-256 (Anfang): \texttt{ad45198c44b84365}.}\par
\subsection*{S17 — Quantisierte Quellenrückwirkung}\label{s:17}
{\small\path{UR.COMPILER.ROOT_SOURCE_BACKREACTION.17}\par}
Vertrag: \path{TFPT/experiments/theory-contracts/compiler-root-source-backreaction-20260919/contract_index.json}\par{\footnotesize SHA-256 (Anfang): \texttt{4df71c5f267eaf0f}.}\par
Herleitung: \path{TFPT/experiments/theory-contracts/compiler-root-source-backreaction-20260919/PROOF.txt}\par{\footnotesize SHA-256 (Anfang): \texttt{f61bf749daf63c10}.}\par
\subsection*{S18 — Phasentreue Lifts und Readouts}\label{s:18}
{\small\path{UR.COMPILER.ROOT_PHASE_LIFT.18}\par}
Vertrag: \path{TFPT/experiments/theory-contracts/compiler-root-phase-lift-20260919/contract_index.json}\par{\footnotesize SHA-256 (Anfang): \texttt{afb3140619c42796}.}\par
Herleitung: \path{TFPT/experiments/theory-contracts/compiler-root-phase-lift-20260919/PROOF.txt}\par{\footnotesize SHA-256 (Anfang): \texttt{d7465bd8ee3514d2}.}\par
\subsection*{S19 — Überlappende Quellen und starke Phasenkopplung}\label{s:19}
{\small\path{UR.COMPILER.OVERLAP_PHASE_LOCK.19}\par}
Vertrag: \path{TFPT/experiments/theory-contracts/compiler-overlap-phase-lock-20260919/contract_index.json}\par{\footnotesize SHA-256 (Anfang): \texttt{c43c5ee8545fd73a}.}\par
Herleitung: \path{TFPT/experiments/theory-contracts/compiler-overlap-phase-lock-20260919/PROOF.txt}\par{\footnotesize SHA-256 (Anfang): \texttt{e24e640c3c341e2f}.}\par
\subsection*{S20 — Ketten, räumliche Antwort und Stromvervollständigung}\label{s:20}
{\small\path{UR.COMPILER.SPATIAL_RESPONSE.20}\par}
Vertrag: \path{TFPT/experiments/theory-contracts/compiler-spatial-response-20260919/contract_index.json}\par{\footnotesize SHA-256 (Anfang): \texttt{a4913e3504199393}.}\par
Herleitung: \path{TFPT/experiments/theory-contracts/compiler-spatial-response-20260919/PROOF.txt}\par{\footnotesize SHA-256 (Anfang): \texttt{b200ef9a3d082cf3}.}\par
\subsection*{S21 — Native Stromprodukte}\label{s:21}
{\small\path{UR.COMPILER.CURRENT_PRODUCT.21}\par}
Vertrag: \path{TFPT/experiments/theory-contracts/compiler-current-product-20260919/contract_index.json}\par{\footnotesize SHA-256 (Anfang): \texttt{1a2843f7bf358abb}.}\par
Herleitung: \path{TFPT/experiments/theory-contracts/compiler-current-product-20260919/PROOF.txt}\par{\footnotesize SHA-256 (Anfang): \texttt{940131acfceba193}.}\par
\subsection*{S22 — Clock-erzwungene E8-Stromschließung}\label{s:22}
{\small\path{UR.COMPILER.CLOCK_CURRENT_CLOSURE.22}\par}
Vertrag: \path{TFPT/experiments/theory-contracts/compiler-clock-current-closure-20260919/contract_index.json}\par{\footnotesize SHA-256 (Anfang): \texttt{f06dd6132ad73df2}.}\par
Herleitung: \path{TFPT/experiments/theory-contracts/compiler-clock-current-closure-20260919/PROOF.txt}\par{\footnotesize SHA-256 (Anfang): \texttt{786817abb2707ee3}.}\par
\subsection*{S23 — Audit der Paar- und Mischsektoren}\label{s:23}
{\small\path{UR.COMPILER.PAIR_SECTOR_AUDIT.23}\par}
Vertrag: \path{TFPT/experiments/theory-contracts/compiler-pair-sector-audit-20260919/contract_index.json}\par{\footnotesize SHA-256 (Anfang): \texttt{49845a7badcaac7c}.}\par
Herleitung: \path{TFPT/experiments/theory-contracts/compiler-pair-sector-audit-20260919/PROOF.txt}\par{\footnotesize SHA-256 (Anfang): \texttt{b63d5a7a97c331f0}.}\par
\subsection*{S24 — Bindungsantwort und damalige Vollraumannahmen}\label{s:24}
{\small\path{UR.COMPILER.BOND_RESPONSE.24}\par}
Vertrag: \path{TFPT/experiments/theory-contracts/compiler-bond-response-20260919/contract_index.json}\par{\footnotesize SHA-256 (Anfang): \texttt{93c749362eda3d17}.}\par
Herleitung: \path{TFPT/experiments/theory-contracts/compiler-bond-response-20260919/PROOF.txt}\par{\footnotesize SHA-256 (Anfang): \texttt{67e3d788c0bfe616}.}\par
\subsection*{S25 — Domänenwandkern}\label{s:25}
{\small\path{UR.COMPILER.DOMAIN_WALL_CORE.25}\par}
Vertrag: \path{TFPT/experiments/theory-contracts/compiler-domain-wall-core-20260919/contract_index.json}\par{\footnotesize SHA-256 (Anfang): \texttt{35276acd6840b8a3}.}\par
Herleitung: \path{TFPT/experiments/theory-contracts/compiler-domain-wall-core-20260919/PROOF.txt}\par{\footnotesize SHA-256 (Anfang): \texttt{fe9c66ef6821ff47}.}\par
\subsection*{S26 — Relaxierter Kern und Zwei-Quellen-Vollsatz}\label{s:26}
{\small\path{UR.COMPILER.RELAXED_WALL.26}\par}
Vertrag: \path{TFPT/experiments/theory-contracts/compiler-relaxed-wall-20260919/contract_index.json}\par{\footnotesize SHA-256 (Anfang): \texttt{479ba50d2289f2ca}.}\par
Herleitung: \path{TFPT/experiments/theory-contracts/compiler-relaxed-wall-20260919/PROOF.txt}\par{\footnotesize SHA-256 (Anfang): \texttt{b4930e179bc5658c}.}\par
\subsection*{S27 — Relationale Inzidenz}\label{s:27}
{\small\path{UR.COMPILER.RELATIONAL_INCIDENCE.27}\par}
Vertrag: \path{TFPT/experiments/theory-contracts/compiler-relational-incidence-20260920/contract_index.json}\par{\footnotesize SHA-256 (Anfang): \texttt{40c82156b0826b7f}.}\par
Herleitung: \path{TFPT/experiments/theory-contracts/compiler-relational-incidence-20260920/PROOF.txt}\par{\footnotesize SHA-256 (Anfang): \texttt{d86af76b009c6126}.}\par
\subsection*{S28 — Vier Follow-ups und Vier-Quellen-Vollsatz}\label{s:28}
{\small\path{UR.COMPILER.FOUR_FOLLOWUPS.28}\par}
Vertrag: \path{TFPT/experiments/theory-contracts/compiler-four-followups-20260920/contract_index.json}\par{\footnotesize SHA-256 (Anfang): \texttt{f4937cd581079dd4}.}\par
Herleitung: \path{TFPT/experiments/theory-contracts/compiler-four-followups-20260920/PROOF.txt}\par{\footnotesize SHA-256 (Anfang): \texttt{ee0775ff2dbd37ac}.}\par
Herleitung: \path{TFPT/experiments/theory-contracts/compiler-four-followups-20260920/origin/PROOF.txt}\par{\footnotesize SHA-256 (Anfang): \texttt{82bf4dce1ff3a475}.}\par
Herleitung: \path{TFPT/experiments/theory-contracts/compiler-four-followups-20260920/vacuum/PROOF.txt}\par{\footnotesize SHA-256 (Anfang): \texttt{8ba73582b414db25}.}\par
Herleitung: \path{TFPT/experiments/theory-contracts/compiler-four-followups-20260920/large_chain/PROOF.txt}\par{\footnotesize SHA-256 (Anfang): \texttt{fde2c9f5d566a185}.}\par
\subsection*{S29 — Kritische Brücke, Rahmen und Fermi-Patch}\label{s:29}
{\small\path{UR.COMPILER.CRITICAL_BRIDGE.29}\par}
Vertrag: \path{TFPT/experiments/theory-contracts/compiler-critical-bridge-20260920/contract_index.json}\par{\footnotesize SHA-256 (Anfang): \texttt{ae730a091abc621a}.}\par
Herleitung: \path{TFPT/experiments/theory-contracts/compiler-critical-bridge-20260920/PROOF.txt}\par{\footnotesize SHA-256 (Anfang): \texttt{cede7d0ba0b115e1}.}\par
Herleitung: \path{TFPT/experiments/theory-contracts/compiler-critical-bridge-20260920/energy/PROOF.txt}\par{\footnotesize SHA-256 (Anfang): \texttt{9a27358119b20cce}.}\par
Herleitung: \path{TFPT/experiments/theory-contracts/compiler-critical-bridge-20260920/symmetry/PROOF.txt}\par{\footnotesize SHA-256 (Anfang): \texttt{87990f749b795ee6}.}\par
Herleitung: \path{TFPT/experiments/theory-contracts/compiler-critical-bridge-20260920/frames/PROOF.txt}\par{\footnotesize SHA-256 (Anfang): \texttt{1b7f093501662183}.}\par
Herleitung: \path{TFPT/experiments/theory-contracts/compiler-critical-bridge-20260920/origin/PROOF.txt}\par{\footnotesize SHA-256 (Anfang): \texttt{f9a1af5f48253a39}.}\par
\section{Frühe und übergreifende Quellen}
\subsection*{SCTX — Übergeordnete TFPT-Architektur}\label{s:CTX}
Orientierungskarte des bestehenden Korpus, ausdrücklich keine neue Beweisquelle. Der Bericht nutzt die dort beschriebenen Architektur- und T1--T8-Verträge als Maßstab.
Quelle: \path{/Users/stefanhamann/.codex/context/tfpt-gesamtkontext.md}\par{\footnotesize SHA-256 (Anfang): \texttt{662be98c30cd492b}.}\par
\subsection*{SLOG — Primärlog und retrospektive Readout-Auswertung}\label{s:LOG}
Beobachtete Telemetrie, unvollständiger Lauf; keine unabhängige physikalische Validierung der erzeugten Wörter.
Quelle: \path{/Users/stefanhamann/Documents/run_20260919_061343.log}\par{\footnotesize SHA-256 (Anfang): \texttt{626b8365606b397a}.}\par
Quelle: \path{SESSION/outputs/readout-analyse-20260919/Auswertung.txt}\par{\footnotesize SHA-256 (Anfang): \texttt{6aadf68ed6ee45d6}.}\par
Quelle: \path{SESSION/outputs/readout-analyse-20260919/manifest.json}\par{\footnotesize SHA-256 (Anfang): \texttt{d913f4ea108935af}.}\par
Quelle: \path{SESSION/outputs/readout-analyse-20260919/Befunde.json}\par{\footnotesize SHA-256 (Anfang): \texttt{fb649bd50831c908}.}\par
\subsection*{SSTATE — Zustand-Dynamik-Injektivität}\label{s:STATE}
Quelle: \path{TFPT/experiments/theory-contracts/seam-state-dynamics-injectivity-20260918/README.md}\par{\footnotesize SHA-256 (Anfang): \texttt{9b3d02af286051e4}.}\par
Quelle: \path{TFPT/experiments/theory-contracts/seam-state-dynamics-injectivity-20260918/contract_index.json}\par{\footnotesize SHA-256 (Anfang): \texttt{35966a130f0bbd8b}.}\par
Quelle: \path{TFPT/experiments/theory-contracts/seam-state-dynamics-injectivity-20260918/validation.json}\par{\footnotesize SHA-256 (Anfang): \texttt{9f53d4b3f366d047}.}\par
\subsection*{STRANSFER — Auswahl und Nichtauswahl des kohärenten Transfers}\label{s:TRANSFER}
Quelle: \path{SESSION/outputs/transfer-ursprung/transfer_selection_proof.txt}\par{\footnotesize SHA-256 (Anfang): \texttt{6f94273480880627}.}\par
\subsection*{SPHASE — Native Phasenfaser und Holonomie}\label{s:PHASE}
Quelle: \path{SESSION/outputs/native-phasen/Verbindung-und-Phasen.md}\par{\footnotesize SHA-256 (Anfang): \texttt{b18c3af611c08d6e}.}\par
Quelle: \path{SESSION/outputs/native-phasen/Phasenfaser-Beweis.md}\par{\footnotesize SHA-256 (Anfang): \texttt{041129f0be57cfe0}.}\par
Quelle: \path{SESSION/outputs/native-phasen/Holonomie-Beweis.md}\par{\footnotesize SHA-256 (Anfang): \texttt{d33cd0c80658ea3f}.}\par
Quelle: \path{SESSION/outputs/native-phasen/time_one_certificate.json}\par{\footnotesize SHA-256 (Anfang): \texttt{e276cd68049478e7}.}\par
\subsection*{SINCLUSIVE — Inklusive Dreierauslesung und allgemeine Relaxationsschranke}\label{s:INCLUSIVE}
Quelle: \path{SESSION/outputs/dreier-auslesung/Inklusiver-Auslesebeweis.md}\par{\footnotesize SHA-256 (Anfang): \texttt{ed1c5c0b421a7eb1}.}\par
Quelle: \path{SESSION/outputs/dreier-auslesung/Allgemeine-Lueckenschranke.md}\par{\footnotesize SHA-256 (Anfang): \texttt{eb55cf5d9af8d55d}.}\par
\subsection*{SBOUNDARY — Diskrete Rand-Clock-Verbindung}\label{s:BOUNDARY}
Quelle: \path{SESSION/outputs/rand-clock/Rand-und-diskrete-Clock.md}\par{\footnotesize SHA-256 (Anfang): \texttt{6f9aca8838176d19}.}\par
Quelle: \path{SESSION/outputs/rand-clock/certificate.json}\par{\footnotesize SHA-256 (Anfang): \texttt{f56760de8f1b4f16}.}\par
Quelle: \path{SESSION/outputs/rand-clock/continuous_duals.json}\par{\footnotesize SHA-256 (Anfang): \texttt{0ef66b7c9a7b2acf}.}\par
\subsection*{SCOMMON — Gemeinsamer Clockträger und geladene Zielalgebra}\label{s:COMMON}
Quelle: \path{SESSION/outputs/gemeinsamer-ursprung/Universalraum-gemeinsamer-Ursprung.txt}\par{\footnotesize SHA-256 (Anfang): \texttt{db430544be97ba41}.}\par
\section{Dokumentationsaudit und ergänzende Belege}
Die beiden Ergebnissynthesen wurden anhand der jeweiligen Vertragsindizes und Beweisdateien angefertigt. Sie sind redaktionelle Hilfen, keine zusätzlichen unabhängigen Beweise. Die numerische Antworttabelle in Kapitel 5 bewahrt ausdrücklich den Status des eingereichten Sitzungsberichts; sie wird nicht als neue unabhängige Vollrechnung ausgegeben.
Dokumentationsbeleg: \path{SESSION/work/session-documentation-20260920/early_results.txt}\par{\footnotesize SHA-256 (Anfang): \texttt{67a4af8350b47696}.}\par
Dokumentationsbeleg: \path{SESSION/work/session-documentation-20260920/late_results.txt}\par{\footnotesize SHA-256 (Anfang): \texttt{b5149b471c273664}.}\par
Dokumentationsbeleg: \path{SESSION/work/session-documentation-20260920/initial_context.txt}\par{\footnotesize SHA-256 (Anfang): \texttt{505ebbe6be3111c4}.}\par
Dokumentationsbeleg: \path{SESSION/work/session-documentation-20260920/repo_audit.log}\par{\footnotesize SHA-256 (Anfang): \texttt{23161a856a52eefa}.}\par

\renewcommand{\bibname}{Wissenschaftlicher Methodenrahmen}\setbibpreamble{
Die folgenden Primärarbeiten dienen der Einordnung verwendeter Methoden. Die konkreten Quellenidentitäten, endlichen Energieschranken und Ausschlüsse dieses Dossiers werden von den jeweiligen lokalen Vertragsbeweisen getragen. Keine der Literaturreferenzen beweist den fehlenden gemeinsamen TFPT-Ursprung.}

\begin{thebibliography}{9}
\bibitem{sw} Sergey Bravyi, David DiVincenzo und Daniel Loss.
\emph{Schrieffer--Wolff transformation for quantum many-body systems.}
Annals of Physics 326 (2011), 2793--2826.
\url{https://arxiv.org/abs/1105.0675}.
Methodischer Rahmen kontrollierter effektiver Hamiltonoperatoren, nicht Rechtfertigung einer unkontrollierten Projektion.

\bibitem{osborne} Tobias J. Osborne und Alexander Stottmeister.
\emph{Conformal field theory from lattice fermions.}
Communications in Mathematical Physics (2022), arXiv-Version 3.
\url{https://arxiv.org/abs/2107.13834v3}.
Rigorer Vergleichsrahmen für Gitterapproximationen freier Fermion- und WZW-Theorien. Die Identifikation der hiesigen Rohquelle bleibt separat zu beweisen.

\bibitem{coxeter} Ruwen Hollenbach und Patrick Wegener.
\emph{The centralizer of a Coxeter element.} arXiv-Version 2 (2020).
\url{https://arxiv.org/abs/1912.01529v2}.
Hintergrund zur Zentralisatorstruktur. Der weiter gehende konkret verwendete integrale Rahmenausschluss ist in S29 mit seinen eigenen Voraussetzungen dokumentiert.
\end{thebibliography}
\end{document}
